% Separate continuation after HD-F1-S01--S08.
% GPT-6 Astra Pro; /root/fixed_surface; 2026-11-22 UTC.
% Effort not exposed. This is a proof fragment, not a proof-assistant certificate.

\section{HD-F1-S09: the original canonical norm surface}
\label{sec:f1-original-norm-surface}

Let $\pi:Z\to C$ be a connected etale cyclic cover of degree three,
where $C$ is a smooth nonhyperelliptic curve of genus three. Write
$\sigma$ for a deck generator, $X=\operatorname{Sym}^4 Z$, and
\[
 Q=|K_C|\simeq\mathbb P^2,\qquad
 T=(\pi^{(4)})^{-1}Q,\qquad
 a:X\longrightarrow\operatorname{Pic}^4 Z,\qquad
 F=\operatorname{Nm}_\pi\circ a.
 \tag{F1S9.1}
\]
The scheme $T$ is precisely $F^{-1}(K_C)$. These are the actual
symmetric-product and norm morphisms in Schoen's unramified construction.
No identification with the surface of the 2007 smoothing construction
is used in this section.

\paragraph{Smoothness, including repeated points.}
Write an effective divisor as $D=\sum_z m_z z$ and set
\[
 E_{\max}(D)=\sum_{x\in C}\max_{z\in\pi^{-1}(x)}m_z\,x.
\]
The tangent-cotangent description of the Abel morphism identifies the
dual of $dF_D$ with evaluation
\[
 H^0(C,K_C)\xrightarrow{\pi^*}H^0(Z,K_Z)
          \longrightarrow H^0(D,K_Z|_D).
 \tag{F1S9.2}
\]
The equality $K_Z=\pi^*K_C$ uses etaleness. In a local parameter at $x$,
vanishing on the length-$m_z$ subscheme at each lift $z$ requires exactly
vanishing of order $\max_zm_z$ at $x$. Consequently the kernel of
(F1S9.2) is exactly
\[
                         H^0(C,K_C-E_{\max}(D)).
 \tag{F1S9.3}
\]
For $D\in T$, one has $E_{\max}(D)\leq\pi_*D\in|K_C|$ and
$\deg E_{\max}(D)\geq\lceil4/3\rceil=2$. The canonical section of
$\pi_*D$ belongs to (F1S9.3), so this space is nonzero. Its line bundle
has degree at most two. Two independent sections of such a bundle,
after removing a base divisor, would give a map of degree at most two
to $\mathbb P^1$, contrary to nonhyperellipticity and positive genus.
Thus (F1S9.3) has dimension one, and $dF_D$ has rank two at every $D\in T$.

The map $\pi^{(4)}$ is finite flat of degree $3^4=81$: it is finite
between smooth fourfolds, so the Cohen--Macaulay flatness criterion
applies. Its inverse image of $Q$ has pure dimension two. The tangent
space of its scheme fibre $F^{-1}(K_C)$ has dimension $4-2=2$.
Every local ring of $T$ is therefore regular of dimension two.
In particular $T$ is smooth and reduced, without any reduced-support
restriction on the divisors used in the argument.

Let $p:T\to Q$ be the finite morphism. The line (F1S9.3) is the line of
the section defining $p(D)$, because that section already spans it.
Thus the differential of $F$ gives the actual exact conormal sequence
\[
 0\longrightarrow p^*\mathcal O_Q(-1)
  \longrightarrow H^0(C,K_C)\otimes\mathcal O_T
  \longrightarrow N^\vee_{T/X}\longrightarrow0.
 \tag{F1S9.4}
\]
The quotient is $p^*T_{\mathbb P^2}(-1)$ by the Euler sequence. Hence
$N_{T/X}=p^*\Omega_{\mathbb P^2}^1(1)$.

\paragraph{The cycle class and Chern numbers, with normalizations.}
Write $\eta$ for the point-divisor class on $X$, write
$\vartheta=a^*[\Theta_Z]$, and write
$\nu=F^*[\Theta_C]$. The hyperplane class $h=p^*c_1\mathcal O_Q(1)$
is $3\eta|_T$ in cohomology: the inverse image of the point divisor
on $\operatorname{Sym}^4C$ is the sum of the three point divisors of
its lifts. Equation (F1S9.4) consequently gives
\[
 c_1(N_{T/X})=-3\eta|_T,\qquad c_2(N_{T/X})=9\eta^2|_T.
 \tag{F1S9.5}
\]

Here is a derivation of the symmetric-product Chern formulas used below.
For a genus-$g$ curve and its $d$th symmetric power, let $\mathcal D$
be the universal divisor. Its class is
$\delta=d[\mathrm{pt}]\otimes1+1\otimes\eta+\gamma$, where
$\gamma^2=-2[\mathrm{pt}]\otimes\vartheta$ and products of $\gamma$
of degree at least three vanish along the curve factor. The tangent
bundle is the pushforward of $\mathcal O_{\mathcal D}(\mathcal D)$.
The sequence for this divisor and Grothendieck--Riemann--Roch give
\[
 \operatorname{ch}(T\operatorname{Sym}^d Z)
       =(d+1-g-\vartheta)e^\eta+(g-1).
 \tag{F1S9.6}
\]
The identity for $\gamma^2$ follows by using dual symplectic bases of
$H^1(Z)$ and its Jacobian in the Poincare class. In particular, for
$g=7,d=4$,
\[
 c_1(T_X)=-2\eta-\vartheta,\qquad
 c_2(T_X)=3\eta^2+3\eta\vartheta+\vartheta^2/2.
 \tag{F1S9.7}
\]

On $\operatorname{Sym}^4C$ use the classes $x,\theta$. The canonical
plane is the rank-two degeneracy locus of the cotangent differential
$H^0(C,K_C)\otimes\mathcal O\to\Omega^1_{\operatorname{Sym}^4C}$.
This determinantal scheme is reduced with its expected codimension two.
For completeness, near $K_C$ a complex representing the cohomology of
the universal degree-four line bundle has a map of ranks three to one.
Its first derivative at $K_C$ is the perfect Serre pairing
$H^1(\mathcal O_C)\otimes H^0(K_C)\to H^1(K_C)$.
After analytic coordinates this is the row $(u_1,u_2,u_3)$.
The local projectivized kernel is
$u_1s_1+u_2s_2+u_3s_3=0$ in a product with $\mathbb P^2$.
In the chart $s_1=1$, the minors of its cotangent differential
generate $(u_2,u_3)$; this is exactly the reduced canonical plane.
The Porteous formula and (F1S9.6), now for $g=3,d=4$, give
\[
 [Q]=c_2(\Omega^1_{\operatorname{Sym}^4C})
       =x^2-x\theta+\theta^2/2.
\]
Finite flat pullback therefore yields
\[
                         [T]=9\eta^2-3\eta\nu+\nu^2/2.
 \tag{F1S9.8}
\]

The symmetric endomorphism belonging to $\operatorname{Nm}_\pi^*
\Theta_C$ relative to $\Theta_Z$ is $\pi^*\operatorname{Nm}_\pi$.
It has eigenvalue three on $\pi^*H^1(C)$ and zero on the Prym part.
Thus its seven analytic eigenvalues are $3,3,3,0,0,0,0$.
The mixed intersection formula on $\operatorname{Sym}^4 Z$ is
\[
 \int_X\eta^{4-k}(s\vartheta+t\nu)^k
   = k![z^k](1+(s+3t)z)^3(1+sz)^4.
 \tag{F1S9.9}
\]
It follows either by expanding the pulled-back forms on the ordered
product $Z^4$ and dividing by $4!$, or by Poincare's formula for the
Abel image of $\operatorname{Sym}^k Z$ after intersecting the remaining
point divisors. In particular, the terms needed in (F1S9.8) are
\[
\begin{array}{c|rrrrrrrrr}
 \text{monomial}&\eta^4&\eta^3\vartheta&\eta^3\nu&
 \eta^2\vartheta^2&\eta^2\vartheta\nu&\eta^2\nu^2&
 \eta\vartheta^2\nu&\eta\vartheta\nu^2&\vartheta^2\nu^2\\ \hline
 \text{integral}&1&7&9&42&54&54&270&270&1080.
\end{array}
\]
Therefore
\[
 \int_T\eta^2=9,\qquad
 \int_T\eta\vartheta=36,\qquad
 \int_T\vartheta^2=108.
 \tag{F1S9.10}
\]
Adjunction and division of the total Chern classes in (F1S9.5)--(F1S9.7)
now give
\[
 K_T=(\vartheta-\eta)|_T,\qquad
 c_2(T_T)=(\vartheta^2/2-3\eta^2)|_T.
 \tag{F1S9.11}
\]
The total numerical invariants are
\[
                         K_T^2=45,\quad c_2(T)=27,\quad
                         \chi(\mathcal O_T)=6.
 \tag{F1S9.12}
\]

\paragraph{Components and their actual embedded images.}
The complete independent monodromy and contraction proof is
HD-F1-T07 below.
Its monodromy is exactly
$\{(a_1,a_2,a_3,a_4)\in\mathbb F_3^4:\sum a_i=0\}\rtimes\mathfrak S_4$:
the lifted point-line incidence surface is a $\mathbb P^1$-bundle over
$Z$, so its twelve-sheet monodromy is transitive; ordinary tangent
meridians are colored transpositions; a spanning-tree change of labels
gives ordinary $\mathfrak S_4$; transitivity then supplies a nonzero
color, whose conjugates generate every difference $e_i-e_j$.
Consequently $T$ has exactly three connected components $T_j$, of
degree $27$ over $Q$, and $\sigma$ permutes them cyclically.
All classes in (F1S9.10)--(F1S9.12) are deck invariant, so each component
has
\[
 \eta^2=3,\quad\eta\vartheta=12,\quad\vartheta^2=36,
 \qquad K_{T_j}^2=15,\quad c_2(T_j)=9,\quad
 \chi(\mathcal O_{T_j})=2.
 \tag{F1S9.13}
\]

The reduced Abel image $S_j=a(T_j)$ is smooth and embedded in a
connected norm torsor under the Prym $P$, and $T_j\to S_j$ is the
blowup of finitely many distinct points. The exceptional curves are
exactly the complete degree-four pencils $|L|\simeq\mathbb P^1$ with
$\operatorname{Nm}_\pi L=K_C$ in this component. To record why the
conclusion uses actual morphisms: $Z$ cannot be hyperelliptic, since
its free order-three action would preserve a degree-two fibre over
a fixed point of the induced action on $\mathbb P^1$. Hence
$h^0(Z,L)\leq2$ for degree-four $L$, and every positive Abel fibre
is such a $\mathbb P^1$. Equation (F1S9.11) gives $K_T|L|=-1$,
so adjunction gives $|L|^2=-1$. All have the same Hilbert polynomial
for the ample class $\eta$, and have zero Hilbert tangent space
$H^0(\mathcal O_{\mathbb P^1}(-1))$; a finite-type Hilbert scheme
therefore contains only finitely many of them. After contraction,
the differential of the Abel map on the normal bundle of an
exceptional curve factors through
\[
 \mathcal O_{\mathbb P^1}(-1)
       \longrightarrow T_sS_j\otimes\mathcal O_{\mathbb P^1}.
\]
It is nonzero on every fibre, since the Abel differential has kernel
exactly the tangent line to the complete pencil. Every line in
$T_sS_j$ occurs in this tautological inclusion, so the descended
differential is injective. The proper injective descended immersion
is a closed embedding. This verifies the embedding assertion also
at the contracted points.

Let $r$ be the number of exceptional curves in one component; it is
the same on all three components by deck symmetry. These established
maps give
\[
 \deg_\beta S_j=36,\qquad
 K_{S_j}^2=15+r,\qquad c_2(S_j)=9-r,\qquad
 \chi(\mathcal O_{S_j})=2,
 \tag{F1S9.14}
\]
where $\beta=\Theta_Z|_P$ has type $(1,1,3,3)$.
HD-F1-S10 determines $r$ and the conormal map for every such
nonhyperelliptic cover, using these actual images.

\subsection{HD-F1-S11: the hyperelliptic singularity test}
\label{sec:f1-hyperelliptic-norm-singularity}
The nonhyperellipticity hypothesis cannot be dropped from the
smoothness assertion for $T$. On a hyperelliptic genus-three curve,
choose a Weierstrass point $p$; then $K_C\sim4p$. Write
$\pi^{-1}(p)=\{z_1,z_2,z_3\}$ and take
$D=2z_1+z_2+z_3$. This lies in the canonical norm fibre, but
$E_{\max}(D)=2p$ and
$h^0(K_C-2p)=h^0(2p)=2$. Thus the norm differential has rank one
and the tangent space of its two-dimensional fibre has dimension
three. The fibre is singular at $D$. This explicit counterexample
proves that nonhyperellipticity is necessary for the uniform
smoothness statement about the symmetric-product fibre in HD-F1-S09.
